\documentclass[11pt]{article}
\usepackage{amsmath,hyperref}
\title{A gate within nominal lidar range need not be observed at the decision cut}
\author{}\date{}
\begin{document}\maketitle

\paragraph{Consumer and scope.}
The return-window benchmark needs a closure observation at the
courier's initial choice, while it is still at tile 16. The
nominal positions at 16 and 42 are $(100,100,0)$ and $(100,75,0)$.
The alternate mouth checked in the earlier study is near
$(100,62.5,0)$. Consider a level robot at the given positions and
a narrow fitted passage model, where a solid box across the tracked
tube makes it unusable:
\[
[99.5,100.5]\times[62.45,62.55]\times[0.2,0.7] .
\]
The geometry is deliberately constructed. It is not a Fuel mesh,
a spawned rock or a collision test for the actual cave or Husky.
The negative comparison concerns one lidar's returns, not all
available modalities or the local controller's actual map.

\paragraph{The source interface.}
The pinned \href{https://github.com/osrf/subt/blob/083c6a8acb0c25aa1a35e34517e628cd3b52ae77/submitted_models/costar_husky_sensor_config_1/model.sdf}
{CoSTAR Husky configuration} puts its 3D front lidar at height
$0.716$ m on a level base link, with a 100 m range, 1,800 horizontal
samples, 16 vertical samples over $\pm0.261799$ radians, and nominal
10 Hz. The selected \href{https://github.com/osrf/subt/blob/083c6a8acb0c25aa1a35e34517e628cd3b52ae77/submitted_models/costar_husky_sensor_config_1/launch/vehicle_topics.launch}
{model-specific bridge} connects the 3D point stream to the
robot-local ROS \texttt{points} topic. It is not a gate-status
feed. The pinned model spawner invokes this file rather than
the generic file for that topic. The latter contains a rate-relay
option; this specific launch does not include that option.
Even the generic relay's callback forwards an argument without
waiting for a sensor result. A service return does not establish
that a new frame saw the changed gate. Other launch overlays and
native transports were not excluded. Check the actual launch,
subscriptions, frame metadata and delivery, not the generic
configuration by name.

The selected \href{https://github.com/gazebosim/gz-rendering/blob/4d0043361a135b7e43e0583516f1b3ddca9f434d/ogre2/src/Ogre2GpuRays.cc}
{Dome renderer branch source} constructs nominal equally spaced
inclination and azimuth directions. The
\href{https://github.com/gazebosim/gz-sensors/blob/b7148d659ea430e93b9b1b14f3b8162386636a58/src/GpuLidarSensor.cc}
{point producer} reads rendered ranges, applies noise and publishes
the point cloud when it has a connection. The replay computes ideal
geometric rays at those specified angles, independent of cubemap
pixels, noise, lighting, exact installed software builds, and
the robot-specific live frame. It does not execute rendering or
the transport stack.

\paragraph{Early range, late distinction.}
For a level CoSTAR sample from 16, the two central rows are
roughly $0.0614$ and $1.3706$ m at the target plane.
Even with a $\pm0.05$ m height error and any base yaw, the highest
possible negatively inclined row at the box is below $0.113$ m,
and the lowest positively inclined row is above $1.319$ m.
For the bound, use the minimum $y$ separation of either box face
and subtract the maximum horizontal displacement of the lidar mount.
The horizontal range in any direction that reaches the box is at
least this separation. The magnitude of every nominal vertical
inclination is at least one half its row spacing. Thus all ideal
rays miss the box, regardless of horizontal angular sampling.
With fixed background and ideal first-hit return values, the
open and blocked constructions give identical lidar returns.
Ranging noise dependent on those returns cannot add a distinction.

At nominal yaw and height at 42, the closer viewpoint has 23
intersecting ideal rays, all in one ring. That viewpoint can
distinguish the ideal inputs if the box is not occluded and the
unchanged background yields different first hits. It is reached
only after the route choice the bit was supposed to inform.
A different pinned \href{https://github.com/osrf/subt/blob/083c6a8acb0c25aa1a35e34517e628cd3b52ae77/submitted_models/marble_husky_sensor_config_1/model.sdf}
{MARBLE Husky configuration} yields ten hits from 16 at nominal
yaw and height. Its geometry, mount, sample count and rate
differ. This is a positive geometry control, not a comparable
resource result or a guarantee over poses. The ray replay checks
660 sampled yaw-height pairs for the CoSTAR negative, alongside
the analytic strip for all admitted level yaws and heights.

\paragraph{Capture and delivery have their own cutoff.}
Even for a distinguishing beam pattern, a frame must capture
the new condition and reach the decider before it branches.
Under an ideal periodic sample law with arbitrary phase, order a
closure at a scheduled sample after that sample. If the period is
$p$ and the bounded consumption delay is $\ell$, all phases
permit a strictly later sample by a decision budget $D$ exactly
when $p+\ell\le D$. In integer time the next strict capture is
$p-(\text{phase}\bmod p)$ from the change; phase zero witnesses
necessity. The timing checks, including 5,020 finite phase-budget
queries, do not turn a nominal rate into a transport guarantee.
Faster scans alone do not change the geometric negative.

The next test must justify which geometry and timing the actual
decider sees, establish the closure or clearance of the relevant
swept passage rather than classify a single return, and price
any different viewpoint, sensor or rate. If the relevant object
can change after capture, certify validity through the movement
interval or keep a timely fallback. A pre-read map, generic
rate reply, new front cloud or its absence is not by itself a
qualified closure bit. No independent mission or native operating
envelope has been spent or advanced by this preflight.
\end{document}
